
Por ello construimos un \textbf{modelo estructural}.

Modelizar significa conservar aquello que domina la respuesta mecánica y eliminar detalles secundarios.

Por ejemplo:
\[
\text{ala real}\longrightarrow\text{viga equivalente}
\]

o
\[
\text{larguero}\longrightarrow\text{barra de directriz recta}.
\]

\begin{ojo}
Un modelo no es una copia exacta de la realidad. Es una \textbf{idealización útil} de la realidad.
\end{ojo}

% ============================================================
\safechapter{Sólido rígido, sólido deformable e hipótesis básicas}
% ============================================================

\safesection{Sólido rígido}
Un sólido rígido es una idealización en la que las distancias entre sus puntos permanecen constantes.

Las fuerzas pueden producir traslación y rotación, pero no cambio de forma.

La estática elemental estudia principalmente el equilibrio de este modelo mediante
\[
\sum F_x=0,\qquad \sum F_y=0,\qquad \sum M=0.
\]

\safesection{Sólido deformable}
Una estructura real se deforma cuando está cargada.

En un sólido elástico ideal, una vez retirada la carga que originó la deformación, el sólido recupera su configuración inicial.

La Teoría de Estructuras necesita, por tanto, añadir al equilibrio:
\[
\text{cinemática}+\text{ley constitutiva}
\]
porque ahora queremos conocer no solo si el cuerpo está en equilibrio, sino también cómo se deforma.

\safesection{Hipótesis de comportamiento lineal}
En este bloque se trabaja fundamentalmente con un modelo de \textbf{elasticidad lineal}.

Las hipótesis esenciales son:
\begin{enumerate}
  \item desplazamientos suficientemente pequeños;
  \item deformaciones suficientemente pequeñas;
  \item comportamiento material elástico y lineal.
\end{enumerate}

Con ellas podemos plantear muchas ecuaciones sobre la geometría no deformada.

\begin{idea}
Hay dos linealidades diferentes:
\[
\text{linealidad geométrica}
\]
asociada a pequeños desplazamientos/deformaciones, y
\[
\text{linealidad material}
\]
asociada a una relación lineal entre tensión y deformación.
\end{idea}

\safesection{Material homogéneo e isótropo}
\safesubsection{Homogeneidad}
Un material homogéneo posee las mismas propiedades en todos sus puntos.
\[
E(\vect{x})=E=\text{cte.}
\]

\safesubsection{Isotropía}
Un material isótropo posee las mismas propiedades mecánicas en todas las direcciones.

\begin{ojo}
\textbf{Homogéneo} e \textbf{isótropo} no significan lo mismo.

Un material puede ser homogéneo y anisótropo, como ocurre aproximadamente en muchas láminas de material compuesto.
\end{ojo}

\safesection{Problema estático e isotermo}
En el modelo básico se considera:
\begin{itemize}
  \item un problema cuasiestático, sin efectos inerciales relevantes;
  \item temperatura constante o ausencia de deformaciones térmicas.
\end{itemize}
Esto permite centrarnos en la respuesta mecánica debida a cargas.

\safesection{Principio de superposición}
Si el problema es lineal, la respuesta producida por varias cargas puede obtenerse sumando las respuestas individuales.

Si
\[
\mathcal{L}(u_1)=F_1,\qquad \mathcal{L}(u_2)=F_2,
\]
y $\mathcal{L}$ es lineal,
\[
\mathcal{L}(u_1+u_2)=F_1+F_2.
\]
Por tanto,
\[
u(F_1+F_2)=u(F_1)+u(F_2).
\]

Lo mismo puede aplicarse a reacciones, esfuerzos, tensiones y deformaciones.

\begin{errorbox}
No puedes utilizar superposición sin comprobar primero que el problema pueda considerarse lineal.
\end{errorbox}

\safesection{Principio de Saint-Venant}
\index{Saint-Venant, principio de}El principio de Saint-Venant permite sustituir una distribución complicada de carga por otra estáticamente equivalente cuando estudiamos puntos suficientemente alejados de la zona donde se aplica.

Lo importante lejos de la zona local es conservar:
\[
\text{resultante de fuerzas}\qquad\text{y}\qquad\text{resultante de momentos}.
